\documentclass[prd,superscriptaddress,floatfix,nopacs,preprintnumbers,nofootinbib]{revtex4-2}

%% Packages

\usepackage[margin=1.0in]{geometry}
\usepackage{mathrsfs}
\usepackage{amssymb,bm}
\usepackage{amsmath}
\usepackage{mathtools}
\usepackage{slashed}
\usepackage{graphics}
\usepackage{graphicx}
\usepackage{dsfont}
\usepackage{longtable}
\usepackage{bbm} 
\usepackage{float}
\usepackage{tcolorbox}
\usepackage{lipsum}
\usepackage[hidelinks]{hyperref}




%% Commands

\newcommand{\Deltat}{\boldsymbol{\Delta}_{\perp}}

\newcommand{\etL}[1]{\boldsymbol{\epsilon}_{\perp #1}}
\newcommand{\etU}[1]{\boldsymbol{\epsilon}_{\perp}^{#1}}

\newcommand{\ltL}[1]{\boldsymbol{l}_{\perp #1}}
\newcommand{\ltU}[1]{\boldsymbol{l}_{\perp}^{#1}}

\newcommand{\vtL}[1]{\boldsymbol{v}_{\perp #1}}
\newcommand{\vtU}[1]{\boldsymbol{v}_{\perp}^{#1}}

\newcommand{\wtL}[1]{\boldsymbol{w}_{\perp #1}}
\newcommand{\wtU}[1]{\boldsymbol{w}_{\perp}^{#1}}

\newcommand{\wtCL}[1]{\boldsymbol{w}'_{\perp #1}}
\newcommand{\wtCU}[1]{\boldsymbol{w}_{\perp}^{' #1}}


\newcommand{\LtL}[1]{\boldsymbol{L}_{\perp #1}}
\newcommand{\LtU}[1]{\boldsymbol{L}_{\perp}^{#1}}

\newcommand{\LtCL}[1]{\boldsymbol{L}'_{\perp #1}}
\newcommand{\LtCU}[1]{\boldsymbol{L}_{\perp}^{' #1}}

\newcommand{\ltoneL}[1]{\boldsymbol{l}_{1\perp #1}}
\newcommand{\ltoneU}[1]{\boldsymbol{l}_{1\perp}^{#1}}

\newcommand{\lttwoL}[1]{\boldsymbol{l}_{2\perp #1}}
\newcommand{\lttwoU}[1]{\boldsymbol{l}_{2\perp}^{#1}}

\newcommand{\qtL}[1]{\boldsymbol{q}_{\perp #1}}
\newcommand{\qtU}[1]{\boldsymbol{q}_{\perp}^{#1}}

\newcommand{\ptL}[1]{\boldsymbol{p}_{\perp #1}}
\newcommand{\ptU}[1]{\boldsymbol{p}_{\perp}^{#1}}

\newcommand{\pgammatL}[1]{\boldsymbol{p}_{\gamma\perp #1}}
\newcommand{\pgammatU}[1]{\boldsymbol{p}_{\gamma\perp}^{#1}}

\newcommand{\ptoneL}[1]{\boldsymbol{p}_{1\perp #1}}
\newcommand{\ptoneU}[1]{\boldsymbol{p}_{1\perp}^{#1}}

\newcommand{\pttwoL}[1]{\boldsymbol{p}_{2\perp #1}}
\newcommand{\pttwoU}[1]{\boldsymbol{p}_{2\perp}^{#1}}

\newcommand{\ktL}[1]{\boldsymbol{k}_{\perp #1}}
\newcommand{\ktU}[1]{\boldsymbol{k}_{\perp}^{#1}}

\newcommand{\pgtL}[1]{\boldsymbol{p}_{g\perp #1}}
\newcommand{\pgtU}[1]{\boldsymbol{p}_{g\perp}^{#1}}

\newcommand{\PtL}[1]{\boldsymbol{P}_{\perp #1}}
\newcommand{\PtU}[1]{\boldsymbol{P}_{\perp}^{#1}}

\newcommand{\gammatL}[1]{\boldsymbol{\gamma}_{\perp #1}}
\newcommand{\gammatU}[1]{\boldsymbol{\gamma}_{\perp}^{#1}}

\newcommand{\GammatLU}[2]{\boldsymbol{\Gamma}_{\perp #1}^{#2}}

\newcommand{\GammatCLU}[2]{\boldsymbol{\overline{\Gamma}}_{\perp #1}^{#2}}

\newcommand{\Hcal}{\mathcal{H}}
\newcommand{\Ocal}{\mathcal{O}}
\newcommand{\Mcal}{\mathcal{M}}
\newcommand{\Fcal}{\mathcal{F}}
\newcommand{\Acal}{\mathcal{A}}
\newcommand{\Ncal}{\mathcal{N}}
\newcommand{\Tcal}{\mathcal{T}}
\newcommand{\Scal}{\mathcal{S}}
\newcommand{\Ccal}{\mathcal{C}}
\newcommand{\Pcal}{\mathcal{P}}
\newcommand{\Rcal}{\mathcal{R}}
\newcommand{\Ical}{\mathcal{I}}
\newcommand{\Jcal}{\mathcal{J}}
\newcommand{\Kcal}{\mathcal{K}}


\newcommand{\vect}[1]{\boldsymbol{#1}_{\perp}}
\newcommand{\kt}{\vect{k}}
\newcommand{\pt}{\vect{p}}
\newcommand{\pgt}{\vect{p_g}}
\newcommand{\pgammat}{\vect{p_\gamma}}
\newcommand{\Pt}{\vect{P}}
\newcommand{\Dt}{\vect{\Delta}}
\newcommand{\qt}{\vect{q}}
\newcommand{\lt}{\vect{l}}
\newcommand{\ellt}{\vect{\ell}}
\newcommand{\ellttwo}{\boldsymbol{\ell_{2\perp}}}
\newcommand{\vt}{\vect{v}}
\newcommand{\at}{\vect{a}}
\newcommand{\bt}{\vect{b}}
\newcommand{\Bt}{\vect{B}}
\newcommand{\Kt}{\vect{K}}
\newcommand{\wt}{\vect{w}}
\newcommand{\wtbar}{\vect{\bar w}}
\newcommand{\xt}{\vect{x}}
\newcommand{\yt}{\vect{y}}
\newcommand{\zt}{\vect{z}}
\newcommand{\ut}{\vect{u}}
\newcommand{\bxt}{\vect{\bar x}}
\newcommand{\byt}{\vect{\bar y}}
\newcommand{\bzt}{\vect{\bar z}}
\newcommand{\Xt}{\vect{X}}
\newcommand{\Yt}{\vect{Y}}
\newcommand{\Zt}{\vect{Z}}
\newcommand{\rt}{\vect{r}}
\newcommand{\rtp}{\vect{r}'}
\newcommand{\Rt}{\vect{R}}
\newcommand{\Lt}{\vect{L}}

\newcommand{\kgt}{\boldsymbol{k_{g\perp}}}

\newcommand{\kht}{\boldsymbol{k_{h\perp}}}
\newcommand{\pht}{\boldsymbol{p_{h\perp}}}

\newcommand{\xtn}{\boldsymbol{x_{n\perp}}}
\newcommand{\xtnn}{\boldsymbol{x_{n+1\perp}}}

\newcommand{\xti}{\boldsymbol{x_{i\perp}}}
\newcommand{\xtii}{\boldsymbol{x_{i+1\perp}}}
\newcommand{\xtim}{\boldsymbol{x_{i-1\perp}}}

\newcommand{\lti}{\boldsymbol{l_{i\perp}}}
\newcommand{\ltii}{\boldsymbol{l_{i+1\perp}}}
\newcommand{\ltim}{\boldsymbol{l_{i-1\perp}}}

\newcommand{\ptk}{\boldsymbol{p_{k\perp}}}
\newcommand{\ptone}{\boldsymbol{p_{1\perp}}}
\newcommand{\pttwo}{\boldsymbol{p_{2\perp}}}
\newcommand{\ptj}{\boldsymbol{p_{j\perp}}}
\newcommand{\ktone}{\boldsymbol{k_{1\perp}}}
\newcommand{\kttwo}{\boldsymbol{k_{2\perp}}}
\newcommand{\ltone}{\boldsymbol{l_{1\perp}}}
\newcommand{\lttwo}{\boldsymbol{l_{2\perp}}}
\newcommand{\ltthre}{\boldsymbol{l_{3\perp}}}
\newcommand{\Ltone}{\boldsymbol{L_{1\perp}}}
\newcommand{\Lttwo}{\boldsymbol{L_{2\perp}}}
\newcommand{\Gammat}{\boldsymbol{\Gamma_{\perp}}}
\newcommand{\Ltthre}{\boldsymbol{L_{3\perp}}}
\newcommand{\Lttwox}{\boldsymbol{L_{2x\perp}}}
\newcommand{\Lttwoy}{\boldsymbol{L_{2y\perp}}}
\newcommand{\Ltthrex}{\boldsymbol{L_{3x\perp}}}
\newcommand{\Ltthrey}{\boldsymbol{L_{3y\perp}}}
\newcommand{\wtone}{\boldsymbol{w_{1\perp}}}
\newcommand{\wttwo}{\boldsymbol{w_{2\perp}}}
\newcommand{\rtone}{\boldsymbol{r_{1\perp}}}
\newcommand{\rttwo}{\boldsymbol{r_{2\perp}}}
\newcommand{\ztone}{\boldsymbol{z_{1\perp}}}
\newcommand{\zttwo}{\boldsymbol{z_{2\perp}}}
\newcommand{\xtone}{\boldsymbol{x_{1\perp}}}
\newcommand{\xttwo}{\boldsymbol{x_{2\perp}}}
\newcommand{\rxyt}{\boldsymbol{r}_{xy}}
\newcommand{\ryxt}{\boldsymbol{r}_{yx}}
\newcommand{\rwyt}{\boldsymbol{r}_{wy}}
\newcommand{\rwxt}{\boldsymbol{r}_{wx}}
\newcommand{\rzxt}{\boldsymbol{r}_{zx}}
\newcommand{\rzyt}{\boldsymbol{r}_{zy}}
\newcommand{\rzytp}{\boldsymbol{r}_{z'y'}}
\newcommand{\rxytp}{\boldsymbol{r}_{x'y'}}
\newcommand{\rxxtp}{\boldsymbol{r}_{xx'}}
\newcommand{\ryytp}{\boldsymbol{r}_{yy'}}
\newcommand{\rxypt}{\boldsymbol{r}_{xy'}}
\newcommand{\rzypt}{\boldsymbol{r}_{zy'}}
\newcommand{\rzxpt}{\boldsymbol{r}_{zx'}}
\newcommand{\rzxtp}{\boldsymbol{r}_{z'x'}}
\newcommand{\rxpyt}{\boldsymbol{r}_{x'y}}
\newcommand{\rwytp}{\boldsymbol{r}_{w'y'}}
\newcommand{\et}{\vect{\epsilon}}

\newcommand{\RtS}{\boldsymbol{R}_{\mathrm SE}}
\newcommand{\RtV}{\boldsymbol{R}_{\mathrm V}}
\newcommand{\RtR}{\boldsymbol{R}_{\mathrm R}}
\newcommand{\RtRp}{\boldsymbol{R}_{\mathrm R'}}
\newcommand{\RtSp}{\boldsymbol{R}_{\mathrm SE'}}
\newcommand{\RtVp}{\boldsymbol{R}_{\mathrm V'}}


\newcommand{\Nc}{N_c}
\newcommand{\der}{\mathrm{d}}
\newcommand{\dijet}{2\textrm{jet}}
\newcommand{\Tr}{\mathrm{Tr}}


\renewcommand{\arraystretch}{1.2}

\newcommand{\FS}[1]{{\color{red} \bf FS: #1}}


\begin{document}
\maketitle



% \begin{align}
%     &\frac{\partial}{\partial \ln(1/x)} \frac{1}{N_c}\langle \Tr[V(\xt)V^\dagger(\yt)] \rangle_x \nonumber\\
%     & =  \int \frac{\der^2 \zt}{2\pi} \Kcal(\zt,\xt,\yt) \left[  \frac{1}{N_c}\langle \Tr[V(\xt)V^\dagger(\zt)] \rangle_x \frac{1}{N_c}\langle \Tr[V(\zt)V^\dagger(\yt)]\rangle_x  - \frac{1}{N_c}\langle \Tr[V(\xt)V^\dagger(\yt)] \rangle_x \right]
% \end{align}
% \begin{align}
%     \Kcal(\zt,\xt,\yt) = \frac{\alpha_s N_c}{\pi} \frac{(\xt-\yt)^2}{(\xt-\zt)^2 (\zt-\yt)^2}
% \end{align}

% Take $\yt \to \infty$:
% \begin{align}
%     &\frac{\partial}{\partial \ln(1/x)} \frac{1}{N_c}\langle \Tr[V(\xt)] \rangle_x \nonumber\\
%     & =  \int \frac{\der^2 \zt}{2\pi} \frac{\alpha_s N_c}{\pi} \frac{1}{(\xt-\zt)^2} \left[  \frac{1}{N_c}\langle \Tr[V(\xt)V^\dagger(\zt)]\rangle_x  \frac{1}{N_c}\langle \Tr[V(\zt)]\rangle_x  - \frac{1}{N_c}\langle \Tr[V(\xt)] \rangle_x \right]
% \end{align}

% \newpage

% The JIMWLK hamiltonian is given by
% \begin{align}
%     H_{\rm{LO}} = \frac{1}{2} \int \der^2 \ut \der^2 \vt \frac{\delta}{\delta A^{+,a}_{cl}(\ut)} \eta^{ab}(\ut,\vt) \frac{\delta}{\delta A^{+,b}_{cl}(\vt)} \label{eq:JIMWLK-Hamiltonian}
% \end{align}
% \begin{align}
%     H_{LO,R} = \frac{1}{2} \int \der^2 \ut \der^2 \vt  \eta^{ab}(\ut,\vt) \frac{\delta}{\delta A^{+,a}_{cl}(\ut)} \frac{\delta}{\delta A^{+,b}_{cl}(\vt)} \label{eq:JIMWLK-Hamiltonian-real}
% \end{align}
% \begin{align}
%     H_{LO,V} = \frac{1}{2} \int \der^2 \ut \der^2 \vt \left(\frac{\delta \eta^{ab}(\ut,\vt)}{\delta A^{+,a}_{cl}(\ut)}  \right) \frac{\delta}{\delta A^{+,b}_{cl}(\vt)} \label{eq:JIMWLK-Hamiltonian-virtual}
% \end{align}
% with 
% \begin{align}
%     \eta^{ab}(\ut,\vt)=\frac{1}{\pi}\int\frac{\der^2\zt}{(2\pi)^2}\frac{(\ut-\zt)(\vt-\zt)}{(\ut-\zt)^2(\vt-\zt)^2}[1+U^\dagger(\ut)U(\vt)-U^\dagger(\ut)U(\zt)-U^\dagger(\zt)U(\vt)]^{ab}
% \end{align}
% The action of the derivative is given by
% \begin{align}
%     \frac{\delta V(\xt)}{\delta A^{+,b}(\vt)} = -ig \delta_{xv} V(\xt) t^b
% \end{align}
% \begin{align}
%     \frac{\delta^2 V(\xt)}{\delta A^{+,a}(\ut) \delta A^{+,b}(\vt)} = -ig \delta_{xv} \frac{\delta V(\xt)}{\delta A^{+,a}(\ut)} t^b = -g^2 \delta_{xv} \delta_{xu} V(\xt) t^a t^b
% \end{align}
% \subsection{Real}
% We have
% \begin{align}
%     H_{LO,R} V(\xt) & = \frac{-g^2}{2} \int \der^2 \ut \der^2 \vt  \eta^{ab}(\ut,\vt) \delta_{xv} \delta_{xu} V(\xt) t^a t^b \\
%     &= \frac{-g^2}{2}  \eta^{ab}(\xt,\xt) V(\xt) t^a t^b
% \end{align}
% Note
% \begin{align}
%     \eta^{ab}(\xt,\xt)=\frac{1}{\pi}\int\frac{\der^2\zt}{(2\pi)^2}\frac{1}{(\xt-\zt)^2}[2-U^\dagger(\xt)U(\zt)-U^\dagger(\zt)U(\xt)]^{ab}
% \end{align}
% \begin{align}
%     H_{LO,R} V(\xt) & = \frac{-g^2}{2}  \frac{1}{\pi}\int\frac{\der^2\zt}{(2\pi)^2}\frac{1}{(\xt-\zt)^2}[2-U^\dagger(\xt)U(\zt)-U^\dagger(\zt)U(\xt)]^{ab} V(\xt) t^a t^b
% \end{align}
% Note
% \begin{align}
%    [U^\dagger(\zt)U(\xt)]^{ab}  t^a t^b & =(U^\dagger)^{ac}(\zt)U^{cb}(\xt) t^a t^b \\
%    & =U^{ca}(\zt)U^{cb}(\xt) t^a t^b \\
%    & = V^\dagger(\zt) t^c V(\zt) V^\dagger(\xt) t^c V(\xt)
% \end{align}
% then
% \begin{align}
%     &H_{LO,R} V(\xt) \nonumber \\
%     & = \frac{-g^2}{2}  \frac{1}{\pi}\int\frac{\der^2\zt}{(2\pi)^2}\frac{1}{(\xt-\zt)^2} V(\xt) ( 2 C_F - V^\dagger(\xt) t^c V(\xt) V^\dagger(\zt) t^c V(\zt) - V^\dagger(\zt) t^c V(\zt) V^\dagger(\xt) t^c V(\xt) )
% \end{align}
% \begin{align}
%     \Tr[V(\xt) V^\dagger(\xt) t^c V(\xt) V^\dagger(\zt) t^c V(\zt)] &= \Tr[t^c V(\xt) V^\dagger(\zt) t^c V(\zt)] \\
%     & = \frac{1}{2}\Tr[V(\zt)] \Tr[ V(\xt) V^\dagger(\zt)] - \frac{1}{2N_c} \Tr[V(\xt)]
% \end{align}
% \begin{align}
%     \Tr[V(\xt) V^\dagger(\zt) t^c V(\zt) V^\dagger(\xt) t^c V(\xt) ] =\frac{1}{2} \Tr[ V(\xt) V(\xt) V^\dagger(\zt)] \Tr[V(\zt) V^\dagger(\xt)] - \frac{1}{2N_c} \Tr[V(\xt)]
% \end{align}
% Thus
% \begin{align}
%     &H_{LO,R} \Tr[V(\xt)] = \frac{-g^2}{2}  \frac{1}{\pi}\int\frac{\der^2\zt}{(2\pi)^2} \frac{1}{(\xt-\zt)^2} \nonumber \\
%     &  \left(N_c \Tr[V(\xt) ] - \frac{1}{2}\Tr[V(\zt)] \Tr[ V(\xt) V^\dagger(\zt)] - \frac{1}{2} \Tr[ V(\xt) V(\xt) V^\dagger(\zt)] \Tr[V(\zt) V^\dagger(\xt)] \right)
% \end{align}

% \subsection{Virtual}
% The virtual part is computed in a similar fashion with the help of the following two identities:
% \begin{align}
%     \frac{1}{2} \int \der^2 \ut  \left(\frac{\delta \eta^{ab}(\ut,\vt)}{\delta A^{+,a}_{cl}(\ut)}  \right) = \frac{ig}{2\pi} \int \frac{\der^2 \zt}{(2\pi)^2} \frac{1}{(\vt-\zt)^2} \Tr\left[T^b U^\dagger(\vt) U(\zt) \right]
% \end{align}

% \begin{align}
%     \Tr\left[T^a U^\dagger(\vt) U(\zt) \right] = \Tr[V^\dagger(\vt) V(\zt)t^a ] \Tr[V(\vt) V^\dagger(\zt)] - (v \leftrightarrow z) \label{eq:Tr-adjoint-to-fundamental}
% \end{align}
% Then
% \begin{align}
%     H_{LO,V} V(\xt) & = \frac{-ig }{2} \int \der^2 \ut \der^2 \vt \left(\frac{\delta \eta^{ab}(\ut,\vt)}{\delta A^{+,a}_{cl}(\ut)}  \right)  \delta_{xv} V(\xt) t^b \\
%     & = \frac{-ig }{2} \int \der^2 \ut \left(\frac{\delta \eta^{ab}(\ut,\xt)}{\delta A^{+,a}_{cl}(\ut)}  \right)  V(\xt) t^b \\
%     & = \frac{g^2}{2\pi} \int \frac{\der^2 \zt}{(2\pi)^2} \frac{1}{(\xt-\zt)^2} \Tr\left[T^b U^\dagger(\xt) U(\zt) \right]  V(\xt) t^b \\
%     & = \frac{g^2}{2\pi} \int \frac{\der^2 \zt}{(2\pi)^2} \frac{1}{(\xt-\zt)^2} (\Tr[V^\dagger(\xt) V(\zt)t^b ] \Tr[V(\xt) V^\dagger(\zt)] - (x \leftrightarrow z) ) V(\xt) t^b
% \end{align}
% Note
% \begin{align}
%     \Tr[V^\dagger(\xt) V(\zt)t^b ] \Tr[ V(\xt) t^b] = \frac{1}{2} \Tr[V(\zt)] -\frac{1}{2N_c} \Tr[V^\dagger(\xt) V(\zt)] \Tr[ V(\xt)] 
% \end{align}
% \begin{align}
%     \Tr[V^\dagger(\zt) V(\xt)t^b ] \Tr[ V(\xt) t^b] = \frac{1}{2} \Tr[V^\dagger(\zt) V(\xt) V(\xt)] -\frac{1}{2N_c} \Tr[V^\dagger(\zt) V(\xt)] \Tr[ V(\xt)] 
% \end{align}
% then
% \begin{align}
%     &H_{LO,V} V(\xt) 
%     = \frac{g^2}{2\pi} \int \frac{\der^2 \zt}{(2\pi)^2} \frac{1}{(\xt-\zt)^2} \nonumber \\
%     & \left(\frac{1}{2}\Tr[V(\zt)] \Tr[V(\xt) V^\dagger(\zt)] - \frac{1}{2}\Tr[V^\dagger(\zt) V(\xt) V(\xt)] \Tr[V(\zt) V^\dagger(\xt)]\right)
% \end{align}
% \subsection{Real + Virtual}
% \begin{align}
%     &H_{LO} \Tr[V(\xt)] = \frac{g^2 N_c }{2\pi}\int\frac{\der^2\zt}{(2\pi)^2} \frac{1}{(\xt-\zt)^2} \left( \frac{1}{N_c }\Tr[V(\zt)] \Tr[ V(\xt) V^\dagger(\zt)] -\Tr[V(\xt) ]  \right)
% \end{align}
% \begin{align}
%     &H_{LO} \frac{1}{N_c } \Tr[V(\xt)] = \frac{\alpha_s N_c }{\pi}\int\frac{\der^2\zt}{2\pi} \frac{1}{(\xt-\zt)^2} \left( \frac{1}{N_c }\Tr[V(\zt)] \frac{1}{N_c } \Tr[ V(\xt) V^\dagger(\zt)] -\frac{1}{N_c } \Tr[V(\xt) ]  \right)
% \end{align}

% \newpage


The differential cross-section for coherent vector meson production:
\begin{align}
    \frac{\der \sigma^{\gamma^* + A \to V + A}}{\der |t|}  = \frac{1}{16\pi} \left|\left\langle \mathcal{A} \right\rangle_\Omega\right|^2
\end{align}
Note
\begin{align}
    \der |t| = \der |\Deltat|^2 = 2 |\Deltat| \der |\Deltat| = \frac{1}{\pi}\der^2 \Deltat
\end{align}
Then we can write the differential cross-section as
\begin{align}
    \frac{\der \sigma^{\gamma^* + A \to V + A}}{\der^2 \Deltat}  = \frac{1}{16\pi^2} \left|\left\langle \mathcal{A} \right\rangle_\Omega\right|^2
\end{align}
\begin{align}
    \sigma^{\gamma^* + A \to V + A}  = \frac{1}{16\pi^2} \int \der^2 \Deltat \left|\left\langle \mathcal{A} \right\rangle_\Omega\right|^2
\end{align}
Recall
\begin{align}
    \mathcal{A} = 2i\int \der^2 {\rt} \der^2 \bt  \frac{\der z}{4\pi} e^{-i \bt \cdot \Deltat} [\Psi_V^* \Psi_\gamma] (Q^2,\rt,z) N_\Omega(\rt,\bt,z,x)
\end{align}
where\footnote{Off-forward phase has been absorbed into the dipole coordinates}
\begin{align}
    N_\Omega(\rt,\bt,z,x) = 1-\frac{1}{N_c} \Tr\left[V(\bt + (1-z)\rt) V^\dagger(\bt - z\rt) \right]
\end{align}
Let us define
\begin{align}
    F(\bt,x) = \int \der^2 {\rt} \der z   [\Psi_V^* \Psi_\gamma] (Q^2,\rt,z) N_\Omega(\rt,\bt,z,x) 
\end{align}
Then
\begin{align}
    \mathcal{A} = i \int \frac{\der^2 \bt}  {2\pi} e^{-i \bt \cdot \Deltat} F(\bt,x)
\end{align}
Then
\begin{align}
    |\langle \mathcal{A} \rangle_{\Omega}|^2 &= \int \frac{\der^2 \bt}  {2\pi} e^{-i \bt \cdot \Deltat} \langle F(\bt,x) \rangle_{\Omega} \int \frac{\der^2 \bt'}  {2\pi} e^{i \bt' \cdot \Deltat} \langle F^{*}(\bt',x) \rangle_{\Omega}
\end{align}
The integral over $\Dt$ is trivial as it only appears on the phase giving a $(2\pi)^2\delta^{(2)}(\bt-\bt')$, we then integrate over $\bt'$
\begin{align}
    \int \der^2 \Dt |\langle \mathcal{A} \rangle_{\Omega}|^2 &= \int \der^2 \Dt 
 \int \frac{\der^2 \bt}  {2\pi} e^{-i \bt \cdot \Deltat} \langle F(\bt,x) \rangle_{\Omega} \int \frac{\der^2 \bt'}  {2\pi} e^{i \bt' \cdot \Deltat} \langle F^{*}(\bt',x) \rangle_{\Omega} \\
 & =  \int \der^2 \bt \langle F(\bt,x) \rangle_{\Omega} \langle F^{*}(\bt,x) \rangle_{\Omega} \nonumber \\
 & = \int \der^2 \bt |\langle F(\bt,x) \rangle_{\Omega}|^2
\end{align}
Finally, we have
\begin{align}
    \sigma^{\gamma^* + A \to V + A}  = \frac{1}{16\pi^2} \int \der^2 \bt |\langle F(\bt,x) \rangle_{\Omega}|^2
\end{align}
where
\begin{align}
    F(\bt,x) = \int \der^2 {\rt} \der z   [\Psi_V^* \Psi_\gamma] (Q^2,\rt,z) N_\Omega(\rt,\bt,z,x) 
\end{align}
Event-by-event $F(\bt,x)$ might depend on the angle of $\bt$ but after average it only depends on $|\bt|$, so for a given $x$, $\langle F(|\bt|,x) \rangle_{\Omega}$ can be stored in a one-dimensional table. $\langle F(|\bt|,x) \rangle_{\Omega}$ should be a smooth function of $|\bt|$, it can be understood as an $x$-dependent "effective" thickness function \footnote{In the dilute limit in which $\langle N_\Omega(\rt,\bt,z,x) 
\rangle_{\Omega} = \langle N_\Omega(\rt,x) \rangle_{\Omega} T(\bt)$, then $\langle F(\bt,x) \rangle_{\Omega} \propto T(\bt)$}.
\begin{align}
    \sigma^{\gamma^* + A \to V + A}  = \frac{1}{8\pi} \int |\bt| \der |\bt| |\langle F(|\bt|,x) \rangle_{\Omega}|^2
\end{align}
Then the total coherent cross-section is computed from a one-dimensional integral over $|\bt|$.

In order to compute the incoherent cross section, one has to store $F(|\bt|,\theta)$ for every configuration and compute
The incoherent cross-section is
\begin{align}
    \sigma^{\gamma^* + A \to V + A^*}  = \frac{1}{16\pi^2} \int \der^2 \bt \Big[\langle |F(\bt,x) |^2 \rangle_{\Omega} -  |\langle F(\bt,x) \rangle_{\Omega} |^2  
 \Big].
\end{align}


\end{document}
